\documentclass[12pt,a4paper]{article}

\usepackage[utf8]{inputenc}
\usepackage[T1]{fontenc}
\usepackage[margin=2.5cm]{geometry}
\usepackage[dvipsnames]{xcolor}
\usepackage{hyperref}
\usepackage{microtype}
\usepackage{parskip}
\usepackage{lmodern}

\hypersetup{colorlinks=true, urlcolor=blue!60!black}

\pagestyle{empty}

\begin{document}

\begin{flushright}
Kundai Farai Sachikonye \\
Germany \\
\href{https://github.com/fullscreen-triangle}{github.com/fullscreen-triangle} \\
\texttt{kundai.sachikonye@bitspark.com} \\[6pt]
6 July 2026
\end{flushright}

\vspace{1em}

Prof.\ Dr.\ Cheng-Jian Xu \\
Clinical Bioinformatics Group \\
Centre for Individualised Infection Medicine (CiiM) \\
Hannover Medical School (MHH) \\
Hannover

\vspace{1.5em}

\textbf{Re: PhD Position (f/d/m) in Clinical Bioinformatics}

\vspace{1em}

Dear Prof.\ Xu,

I am applying for the PhD position in Clinical Bioinformatics in your group.
The project's core problem --- reading immune aging, multimorbidity, and disease
susceptibility off integrated multi-omics and population data, and turning that
reading into predictive models for stratification --- is the exact problem I have
been working on independently for the past year, both as computational
infrastructure and as a formal method for representing disease state across
biological scales. I want to describe that work concretely, because the fit is
specific rather than general.

The central object in my work is a \emph{disease-state vector}. In my paper
\textit{Syndrome: Categorical Resolution of Disease Trajectories}, a patient's
pathological signature is encoded as a vector $\mathbf{D}\in[0,1]^8$ over eight
cellular oscillator classes --- protein folding, enzymatic catalysis, channel
gating, membrane potential, ATP synthesis, genetic expression, calcium
signalling, and circadian rhythm --- each contributing a coherence index
$\eta_i\in[0,1]$. This is directly a multimorbidity representation: distinct
conditions occupy distinct regions of the same low-dimensional space, so
comorbidity, trajectory, and patient stratification become geometric questions
on a shared manifold rather than separate disease labels. Optimal intervention
targets are selected as $\arg\max_i w_i D_i$, a sparse $\ell_1$ linear programme
with a computable a priori side-effect bound, and disease trajectories resolve in
$\mathcal{O}(N\log N)$. This is the kind of predictive, stratifying model your
project aims to build, and it is designed to take its inputs from exactly the data
types you work with.

A companion paper, \textit{On Disease Epistemology in Blind Receiver}, addresses
the diagnostic criterion itself. It proves that template-based diagnosis ---
comparison against a healthy reference --- is epistemologically insufficient, and
derives the unique viable alternative: loop-holonomy self-consistency of the
biochemical circuit. A system is healthy iff the composed constraint map around
every reaction cycle returns the identity ($H_\ell=\mathrm{Id}\ \forall\ell$);
disease is $H_\ell\neq\mathrm{Id}$ for some cycle. The holonomy test reaches
AUC $=1.00$ against $0.896$ for template comparison across 25 validation
experiments. The relevance to immune aging is direct: aging is not a single-marker
event but a gradual loss of circuit-level consistency that individual biomarkers
miss precisely because it is distributed across the network. A criterion that
detects distributed inconsistency is the right instrument for a process that is
itself distributed.

On the immune side specifically, my paper \textit{On the Thermodynamic
Consequences of Categorical Completion on Biological Systems} models adaptive
immunity as categorical filtering rather than molecular complementarity, and
yields an empirical result on real immunopeptidomics data: MHC binding affinity
correlates with parent-protein categorical richness independently of classical
sequence features ($r=0.73$, $p<10^{-12}$, $n=2{,}847$ IEDB peptides for MHC-I),
and adding this term to NetMHCpan improves AUC from 0.68 to 0.81 ($p<0.001$,
DeLong test). This is the immune-function layer your project studies, treated
with the same modelling discipline.

The computational infrastructure behind this is built and validated against
authoritative reference data, and it maps onto the omics layers named in your
posting. \href{https://github.com/fullscreen-triangle/gospel}{gospel} integrates
whole-genome VCF data with RNA-seq expression matrices at $10^6$ variants/second,
validated against 1000~Genomes Phase 3, gnomAD v3.1.2, and UK~Biobank --- this is
genomic and transcriptomic integration at population scale, the core technical
requirement of the role.
\href{https://github.com/fullscreen-triangle/lavoisier}{lavoisier} is a
mass-spectrometry metabolomics pipeline achieving 98.9\% feature-extraction
accuracy against NIST reference libraries, with distributed computing via Ray and
PyTorch models --- the metabolomic omics layer. Both are deployable and
documented, not one-off scripts; I build reproducible, reference-validated tools,
and several are available in the browser without installation
(\href{https://syndrome-seven.vercel.app/tools}{syndrome-seven.vercel.app}).

The disease-vector formulation extends naturally to your epigenetic focus: the
epigenome enters as an additional coherence layer, and a methylation-derived
biological-aging coordinate sits directly alongside the transcriptomic and
genomic ones on the same manifold. The statistical and integration machinery
I have built is layer-agnostic by design, so adding an omics modality is a
question of ingestion and preprocessing rather than of reworking the model.

My academic background is in computational lipidomics at the Technische
Universit\"at M\"unchen, where I trained in the multi-omics statistical methods
that underpin all of this work; R and Python are both standard tools from that
training. I am fully authorised to work in Germany without restriction or
sponsorship (permanent residence), and I am available for the 1 September 2026
start. I am not proposing to import my framework wholesale into your programme
--- that is not the purpose of a doctoral position. I am proposing that a
candidate who has independently derived and empirically validated a
low-dimensional model of disease state, and built the multi-omics integration
infrastructure your project depends on, will contribute to your work on immune
aging and stratification from early on. I would welcome the chance to discuss it
with you directly.

\vspace{1.5em}
Yours sincerely,

\vspace{2.5em}
\textbf{Kundai Farai Sachikonye}

\end{document}
